% !TeX program = xelatex
%==============================================================
%  Python金融数据分析 · 章节幻灯片模板
%--------------------------------------------------------------
%  使用方法：
%    1. 将本文件夹整体复制一份，重命名为 "Chapter2 slides" 等；
%    2. 把本文件重命名为 chapter2.tex（与章节对应）；
%    3. 修改下方【首页信息设置】中的标题、副标题；
%    4. 在正文区按 \section + frame 组织内容，示例见下文；
%    5. 图片放入 figures/ 文件夹，直接 \includegraphics{文件名} 即可；
%    6. 编译命令（需要 TeX Live，用 xelatex 编译两遍）：
%         xelatex chapterN.tex
%         xelatex chapterN.tex
%    7. 参考文献写入 reference.bib，正文用 \citep{key} 引用。
%==============================================================
\documentclass[9pt]{ctexbeamer}

%%%%-----导入宏包-----%%%%
\usepackage{style/shubeamer}   % SHU 主题（配色、帧标题、页脚、Python代码样式）
\usepackage{amsmath,amsfonts,amssymb,bm}
\usepackage{graphicx,hyperref,url}
\usepackage{subcaption}
\usepackage{booktabs}
\usepackage[super,square]{natbib}
\usepackage{wrapfig}          % 文字绕排
\usepackage{calligra}         % 结尾页花体字

%%%%-----设置字体-----%%%%
\setmainfont{CMU Serif}
\setsansfont{CMU Sans Serif}

%%%%-----常用自定义颜色（正文可直接使用）-----%%%%
\definecolor{nvidia}{RGB}{102,156,28}
\definecolor{red}{RGB}{184,13,73}
\definecolor{orange}{RGB}{242,151,36}
\definecolor{darkteal}{RGB}{43,106,108}
\definecolor{darkblue}{RGB}{4,37,58}

%%%%-----每节开始时自动显示当前节目录-----%%%%
\AtBeginSection[]
{
	\begin{frame}<beamer>
		\frametitle{\textbf{目录}}
		\textbf{\tableofcontents[currentsection]}
	\end{frame}
}

%%%%-----设定图片存放路径-----%%%%
\graphicspath{{figures/}}

%%%%-----首页信息设置（★ 每章只需修改这里 ★）-----%%%%
\AtBeginDocument{%
	%%%%----短标题[显示在页脚]{封面标题}
	\title[Python金融数据分析 · 第N章]{\fontsize{16pt}{22pt}\selectfont {Python金融数据分析}}
	%%%%----副标题（本章主题）
	\subtitle{\fontsize{9pt}{14pt}\selectfont \textit{第N章\quad 在此填写本章标题}}
	%%%%----个人信息
	\author[王伟冠]{%
		\begin{tabular}{ll}%
			主讲人： & 王伟冠\tabularnewline%
			院\quad 系： & 经济学院金融系\tabularnewline
		\end{tabular}
	}
	%%%%----机构信息
	\institute[SHU]{上海大学}
	%%%%----日期信息
	\date[\today]{\today}}

%%%%-----数学宏（向量、矩阵、算子等，见 style/macros.tex）-----%%%%
\input{style/macros.tex}


\begin{document}

%% 封面
\begin{frame}
	\titlepage
\end{frame}

%% 目录页
\section*{目录}
\begin{frame}
	\frametitle{\textbf{目录}}
	\textbf{\tableofcontents}
\end{frame}


%==============================================================
%  正文开始：按 \section{节标题} 组织，每页一个 frame
%==============================================================
\section{第一节标题}

\begin{frame}{要点列表页}
	\begin{itemize}
		\item 一级要点：使用 \texttt{itemize} 环境
		\item 支持嵌套：
		\begin{itemize}
			\item 二级要点自动使用亮蓝色短横线
		\end{itemize}
		\item 重点内容可用 \alert{alert 高亮} 或 \textbf{加粗}
	\end{itemize}
\end{frame}

\begin{frame}{双栏对比页}
	\begin{columns}[T,totalwidth=\textwidth]
		\begin{column}{0.48\textwidth}
			\textbf{左栏标题}
			\begin{itemize}
				\item 要点一
				\item 要点二
			\end{itemize}
		\end{column}
		\begin{column}{0.48\textwidth}
			\textbf{右栏标题}
			\begin{itemize}
				\item 要点一
				\item 要点二
			\end{itemize}
		\end{column}
	\end{columns}
\end{frame}

\begin{frame}{Block 盒子页}
	\begin{block}{定义}
		普通 block：适合定义、结论。\quad 例：$\mathbb{E}[X] = \int x f(x)\,\mathrm{d}x$
	\end{block}
	\begin{alertblock}{注意}
		alertblock：适合强调易错点、注意事项。
	\end{alertblock}
	\begin{exampleblock}{示例}
		exampleblock：适合举例说明。
	\end{exampleblock}
\end{frame}

\section{代码与图表}

\begin{frame}[fragile]{Python 代码页}
	% 含代码的 frame 必须加 [fragile]
	% 代码样式已在 style/shubeamer.sty 中统一定义（style=shupython）
	\begin{lstlisting}
import numpy as np
import pandas as pd

# 读取金融数据并计算收益率
df = pd.read_csv("data.csv", index_col=0, parse_dates=True)
rets = np.log(df / df.shift(1)).dropna()
print(rets.describe())
\end{lstlisting}
\end{frame}

\begin{frame}{插图页}
	\begin{figure}
		\centering
		% 图片放在 figures/ 文件夹中，无需写路径
		% \includegraphics[width=0.6\textwidth]{example-image}
		\caption{图标题示例}
	\end{figure}
	% 双图并排可使用 subcaption：
	% \begin{subfigure}{0.45\textwidth} ... \end{subfigure}\hfill
	% \begin{subfigure}{0.45\textwidth} ... \end{subfigure}
\end{frame}

\begin{frame}{表格页}
	\begin{table}
		\centering
		\caption{表标题示例}
		\begin{tabular}{lcc}
			\toprule
			项目 & 数值A & 数值B \\
			\midrule
			行一 & 1.23 & 4.56 \\
			行二 & 7.89 & 0.12 \\
			\bottomrule
		\end{tabular}
	\end{table}
\end{frame}

%==============================================================
%  附录：参考文献 + 结尾页
%==============================================================
\appendix

\begin{frame}[allowframebreaks]{参考文献}
	\bibliographystyle{style/gbt7714-2005}
	\bibliography{reference.bib}
\end{frame}

\begin{frame}
	\begin{center}
		{\Huge\calligra Thanks for your attention!}
		\vspace{1cm}

		{\Huge Q \& A}
	\end{center}
\end{frame}

\end{document}
